% Generated draft: requires mathematical and source review.
\documentclass[11pt]{article}
\usepackage[margin=25mm]{geometry}
\usepackage{fontspec}
\setmainfont{Latin Modern Roman}
\usepackage{amsmath,amssymb,mathtools,unicode-math}
\setmathfont{Latin Modern Math}
\usepackage{enumitem,needspace,xurl,hyperref}
\setlength{\parindent}{0pt}
\setlength{\parskip}{5pt}
\setlength{\emergencystretch}{3em}
\newcommand{\N}{\mathbb{N}}
\newcommand{\Z}{\mathbb{Z}}
\newcommand{\Q}{\mathbb{Q}}
\newcommand{\R}{\mathbb{R}}
\newcommand{\C}{\mathbb{C}}
\newcommand{\F}{\mathbb{F}}
\newcommand{\A}{\mathbb{A}}
\newcommand{\PP}{\mathbb{P}}
\newcommand{\E}{\mathbb{E}}
\newcommand{\eps}{\varepsilon}
\newcommand{\dd}{\,\mathrm{d}}
\begin{document}
\section*{UnsolvedMath ID: 8800092}
\textbf{Source problem code:} catalogue entry
\subsection{Definitions and mathematical statement}
In its customary elliptic-curve setting, Deuring lifting concerns lifting positive-characteristic elliptic data together with a suitable specified endomorphism to characteristic zero, compatibly with reduction. Lifting one specified endomorphism is different from lifting an entire endomorphism ring. The exact version intended by the source is not quoted.

For a lift over a discrete valuation ring \(R\) with fraction field \(K\) and residue field \(k\), an elliptic scheme \(\mathcal E/R\) has generic fiber \(E_K\) and special fiber \(E_k\). Its sections specialize by
\[
\mathcal E(R)\longrightarrow E_k(k).
\]
An analogue for a distinguished arithmetic point would need to specify not only compatible curve/point specialization but preservation of the arithmetic structure defining that point. Merely lifting an arbitrary point is not by itself the requested Deuring-type phenomenon.

Rosen and Silverman ask for such an analogue for Darmon--Heegner points. The record does not define that point construction, the field and reduction data, the endomorphism or other structure to be retained, or the existence/uniqueness/functorial conclusion. Classical Heegner constructions from imaginary-quadratic CM points and real-quadratic Darmon/Stark--Heegner constructions are different settings, with algebraicity itself qualified in some versions. A characteristic-zero algebraic-point lifting statement therefore cannot be silently imposed on every object called a Darmon--Heegner point.

The question is a theorem-development programme, not a supplied universal lifting assertion with complete hypotheses. The research summary reports no verified analogue in the requested sense and mentions broader p-adic/explicit-class-field-theory work without an exact matching theorem.
\subsection{Short English statement}
Develop the intended Deuring-type lifting theorem for Darmon--Heegner points, preserving their defining arithmetic structure. The point construction, reduction data, and precise lifting conclusion are missing.
\subsection{Sources}
\begin{itemize}
\item Local UnsolvedMath registry, permanent ID \(8800092\), exact source code \texttt{catalogue entry}, ``Deuring lifting for Darmon--Heegner points.'' Proposers Michael Rosen and Joseph Silverman.
\item IPAM \emph{Workshop on Number Theory and Cryptography --- Open Problems} (2006), attachment 6669, recorded PDF page 24, printed slide 18; \url{https://www.ipam.ucla.edu/programs/workshops/workshop-i-number-theory-and-cryptography-open-problems/?tab=schedule}. Recorded access 2026-07-29; pdf-and-legacy-ppt visual-audit extraction.
\item The source names Deuring lifting and Darmon--Heegner points, and the OPEN-TRIAGE research summary mentions Darmon's work, but no complete references or definitions of the intended analogue are supplied.
\end{itemize}
\end{document}
